
\chapter{El grupo fundamental: generalidades}
\label{V}
\marginpar{105}

\setcounter{section}{-1}
\section{Introducción}

El presente Seminario es continuación del Seminario de 1960. Nos referiremos
a este último mediante siglas como \Ref{I.9.7}, que significa: Seminario de
Geometría Algebraica, exposición~I, \No9.7. La numeración de las exposiciones
de 1961 continuará la de 1960. Nos referiremos a los \emph{Elementos de
geometría algebraica} de Dieudonné--Grothendieck mediante siglas como
(EGA~I~8.7.3).

La presente exposición resume (con ligeros complementos) las últimas
exposiciones de 1960, que no habían sido redactadas.

Como en 1961, nos limitaremos por regla general a preesquemas localmente
noetherianos, aunque a menudo esta restricción sea inesencial. En la
exposición~\Ref{VI} daremos por supuesta la teoría del descenso fielmente
plano, resumida en el Seminario Bourbaki \No190. Si fuera necesario,
ofreceremos una exposición más detallada en una exposición
ulterior\footnote{\Cf exp\ptbl\ \Ref{VI} y exp\ptbl\ \Ref{VIII}}, una vez que
el lector haya tenido ocasión de convencerse de la utilidad de esta técnica
para la teoría del grupo fundamental.
